% GENERATED FILE. Edit canonical lessons, then run npm run generate:books.
% canonical-source-hash: fnv1a64:7db5230503dab0f3
% canonical-lessons: ZH-C168-zheng, ZH-C168-zha, ZH-C168-dun, ZH-C168-shaoshui, ZH-C168-tang

\chapter{Steamed or Fried}
\label{ch:zh-steamed-or-fried}
\hlchaptermodality{168}

\begin{chapteropening}
\textbf{By the end of this chapter:} \emph{I can talk about steaming, deep-frying, stewing and boiling water, and ask for sugar.}

Complete the last lesson of chapter 168: I can talk about steaming, deep-frying, stewing and boiling water, and ask for sugar.
\end{chapteropening}

\section[zhēng]{zhēng — to steam (food)}

\label{lesson:ZH-C168-zheng}



\begin{quote}
Before the new one: say the Mandarin for to have a power cut, then the Mandarin for a sound.
\end{quote}

\subsection*{You'll want to know: zhēng}
\emph{zhēng} — \textquotedblleft{}to steam (food)\textquotedblright{}.

\begin{quote}
\emph{māma zài zhēng bāozi} — Mum is steaming buns.
\end{quote}

Its characters come later; for now it is a word for the ear.

\begin{cousinweb}
\emph{zhēng} is inherited from Middle Chinese \emph{tsying} (Guangyun, in Baxter's transcription). Under this character the dictionary gives 'fine firewood'; written another way, the same syllable is glossed 'the breath of a fire rising', which is the steaming.
\end{cousinweb}

\begin{grammarlens}[title={What you've built}]
One more word for saying what I do.
\end{grammarlens}

\subsection*{Your turn}
\begin{itemize}
\raggedright
  \item \emph{Say it:} \emph{zhēng}
  \item \emph{Say it:} \emph{zhēng}, once more
  \item \emph{Say it:} \emph{māma zài zhēng bāozi} — Mum is steaming buns.
  \item \emph{Recall:} say the Mandarin for to have a power cut, then the Mandarin for a sound, then say \emph{zhēng} again
\end{itemize}

\begin{tcolorbox}[breakable,colback=teal!4,colframe=teal!35!black,title={Before you move on}]
What does zhēng mean? (\textquotedblleft{}To steam (food)\textquotedblright{}.) Say it once more.
\end{tcolorbox}
\section[zhá]{zhá — to deep-fry}

\label{lesson:ZH-C168-zha}



\begin{quote}
Before the new one: say the Mandarin for a sound, then the Mandarin for to steam.
\end{quote}

\subsection*{You'll want to know: zhá}
\emph{zhá} — \textquotedblleft{}to deep-fry\textquotedblright{}.

\begin{quote}
\emph{jīntiān wǎnshang wǒmen zhá yú} — We're deep-frying fish tonight.
\end{quote}

Its characters come later; for now it is a word for the ear.

\begin{cousinweb}
\emph{zhá} is probably the word the Guangyun writes another way as \emph{dzreap}, glossed 'to scald in hot water' (Guangyun, in Baxter's transcription), which would regularly give zhá; but no dictionary we use joins that word to this character, so its history is not settled.
\end{cousinweb}

\begin{grammarlens}[title={What you've built}]
One more word for saying what I do.
\end{grammarlens}

\subsection*{Your turn}
\begin{itemize}
\raggedright
  \item \emph{Say it:} \emph{zhá}
  \item \emph{Say it:} \emph{zhá}, once more
  \item \emph{Say it:} \emph{jīntiān wǎnshang wǒmen zhá yú} — We're deep-frying fish tonight.
  \item \emph{Recall:} say the Mandarin for a sound, then the Mandarin for to steam, then say \emph{zhá} again
\end{itemize}

\begin{tcolorbox}[breakable,colback=teal!4,colframe=teal!35!black,title={Before you move on}]
What does zhá mean? (\textquotedblleft{}To deep-fry\textquotedblright{}.) Say it once more.
\end{tcolorbox}
\section[dùn]{dùn — to stew}

\label{lesson:ZH-C168-dun}



\begin{quote}
Before the new one: say the Mandarin for to steam, then the Mandarin for to deep-fry.
\end{quote}

\subsection*{You'll want to know: dùn}
\emph{dùn} — \textquotedblleft{}to stew\textquotedblright{}.

\begin{quote}
\emph{tā zài dùn ròu} — He's stewing some meat.
\end{quote}

Its characters come later; for now it is a word for the ear.

\begin{cousinweb}
\emph{dùn} has no older form of this word in any dictionary we use. The Guangyun lists the character as \emph{dwon}, glossed 'a blazing fire' (Guangyun, in Baxter's transcription), but neither that sense nor that sound gives today's word, so its history is not settled.
\end{cousinweb}

\begin{grammarlens}[title={What you've built}]
One more word for saying what I do.
\end{grammarlens}

\subsection*{Your turn}
\begin{itemize}
\raggedright
  \item \emph{Say it:} \emph{dùn}
  \item \emph{Say it:} \emph{dùn}, once more
  \item \emph{Say it:} \emph{tā zài dùn ròu} — He's stewing some meat.
  \item \emph{Recall:} say the Mandarin for to steam, then the Mandarin for to deep-fry, then say \emph{dùn} again
\end{itemize}

\begin{tcolorbox}[breakable,colback=teal!4,colframe=teal!35!black,title={Before you move on}]
What does dùn mean? (\textquotedblleft{}To stew\textquotedblright{}.) Say it once more.
\end{tcolorbox}
\section[shāoshuǐ]{shāoshuǐ — to boil water}

\label{lesson:ZH-C168-shaoshui}



\begin{quote}
Before the new one: say the Mandarin for to deep-fry, then the Mandarin for to stew.
\end{quote}

\subsection*{You'll want to know: shāoshuǐ}
\emph{shāoshuǐ} — \textquotedblleft{}to boil water\textquotedblright{}.

\begin{quote}
\emph{wǒ qù shāoshuǐ nǐmen yào hē chá ma} — I'll put the kettle on: who wants tea?
\end{quote}

Its characters come later; for now it is a word for the ear.

\begin{cousinweb}
\emph{shāoshuǐ} is built from \emph{shāo} 'to burn' and \zh{水} 'water', both inherited from Middle Chinese, \emph{syew} and \emph{sywijX} (Guangyun, in Baxter's transcription), which glosses the first 'fire' and 'to set alight'.
\end{cousinweb}

\begin{grammarlens}[title={What you've built}]
One more word for saying what I do.
\end{grammarlens}

\subsection*{Your turn}
\begin{itemize}
\raggedright
  \item \emph{Say it:} \emph{shāoshuǐ}
  \item \emph{Say it:} \emph{shāoshuǐ}, once more
  \item \emph{Say it:} \emph{wǒ qù shāoshuǐ nǐmen yào hē chá ma} — I'll put the kettle on: who wants tea?
  \item \emph{Recall:} say the Mandarin for to deep-fry, then the Mandarin for to stew, then say \emph{shāoshuǐ} again
\end{itemize}

\begin{tcolorbox}[breakable,colback=teal!4,colframe=teal!35!black,title={Before you move on}]
What does shāoshuǐ mean? (\textquotedblleft{}To boil water\textquotedblright{}.) Say it once more.
\end{tcolorbox}
\section[táng]{táng — sugar; sweets}

\label{lesson:ZH-C168-tang}



\begin{quote}
Before the new one: say the Mandarin for to stew, then the Mandarin for to boil water.
\end{quote}

\subsection*{You'll want to know: táng}
\emph{táng} — \textquotedblleft{}sugar; sweets\textquotedblright{}.

\begin{quote}
\emph{zhège tāng yào fàng yìdiǎnr táng} — This soup needs a little sugar in it.
\end{quote}

Its characters come later; for now it is a word for the ear.

\begin{cousinweb}
\emph{táng} is inherited from Middle Chinese \emph{dang} (Guangyun, in Baxter's transcription), glossed 'malt syrup'.
\end{cousinweb}

\begin{grammarlens}[title={What you've built}]
One more word for saying what I do.
\end{grammarlens}

\subsection*{Your turn}
\begin{itemize}
\raggedright
  \item \emph{Say it:} \emph{táng}
  \item \emph{Say it:} \emph{táng}, once more
  \item \emph{Say it:} \emph{zhège tāng yào fàng yìdiǎnr táng} — This soup needs a little sugar in it.
  \item \emph{Recall:} say the Mandarin for to stew, then the Mandarin for to boil water, then say \emph{táng} again
\end{itemize}

\begin{tcolorbox}[breakable,colback=teal!4,colframe=teal!35!black,title={Before you move on}]
What does táng mean? (\textquotedblleft{}Sugar; sweets\textquotedblright{}.) Say it once more.
\end{tcolorbox}
